Эта дополнительная методичка рассчитана на компьютер, где \textbf{Ubuntu
уже установлена рядом с Windows} и при включении доступен загрузчик GRUB
2. Устанавливать Ubuntu, менять разделы диска или настройки загрузчика
здесь не требуется. После входа в Ubuntu вы создадите учебную папку,
напишете текстовый файл и научитесь находить, копировать и перемещать
файлы. Все примеры выполняйте от имени обычного пользователя.

\hypertarget{ux432ux44bux431ux43eux440-ubuntu-ux432-grub-2}{%
\section{Выбор Ubuntu в GRUB
2}\label{ux432ux44bux431ux43eux440-ubuntu-ux432-grub-2}}

\begin{enumerate}
\def\labelenumi{\arabic{enumi}.}
\tightlist
\item
  Сохраните работу в запущенной системе и перезагрузите компьютер. GRUB
  2 --- это меню выбора операционной системы, которое появляется
  \textbf{до} запуска Ubuntu или Windows.
\item
  Если меню появилось, стрелками выберите пункт \textbf{Ubuntu} (не
  \textbf{Advanced options for Ubuntu} и не \textbf{Windows Boot
  Manager}) и нажмите \textbf{Enter}. Названия и оформление могут
  немного отличаться.
\item
  Дождитесь экрана входа Ubuntu, выберите свою учётную запись и введите
  пароль. При наборе пароля символы могут не отображаться --- это
  нормально. После входа дождитесь рабочего стола.
\end{enumerate}

Если меню GRUB 2 не появилось, но сразу загрузилась Ubuntu, переходите к
следующему разделу: иногда меню скрыто при автоматическом запуске. Если
запустилась Windows либо в меню нет Ubuntu, остановитесь и уточните у
преподавателя или администратора, как выбрать установленную систему
\textbf{на этом компьютере}. Не меняйте разделы диска или параметры
загрузчика ради выполнения этой инструкции.

\hypertarget{ux43eux442ux43aux440ux44bux432ux430ux435ux43c-ux442ux435ux440ux43cux438ux43dux430ux43b}{%
\section{Открываем
терминал}\label{ux43eux442ux43aux440ux44bux432ux430ux435ux43c-ux442ux435ux440ux43cux438ux43dux430ux43b}}

На рабочем столе Ubuntu нажмите \textbf{Ctrl+Alt+T}. Откроется окно
терминала с приглашением ввести команду. Другой способ: откройте обзор
приложений, найдите \textbf{Терминал} (\emph{Terminal}) и запустите его.
Если сочетание клавиш в вашей конфигурации не сработало, используйте
поиск приложений.

\textbf{Терминал} показывает введённые команды и их результат.
\textbf{Командная оболочка} (\emph{shell}, часто Bash) читает команду,
например \texttt{pwd}, и запускает её. В приглашении могут быть имя
пользователя, имя компьютера и знак \texttt{\$} --- \textbf{сам знак
\texttt{\$} из примеров вводить не нужно}. Чтобы выполнить набранную
команду, нажмите \textbf{Enter}. Команда может ничего не напечатать,
если выполнилась успешно.

\hypertarget{ux433ux434ux435-ux43cux44b-ux43dux430ux445ux43eux434ux438ux43cux441ux44f-ux43aux430ux442ux430ux43bux43eux433ux438-ux438-ux43fux443ux442ux438}{%
\section{Где мы находимся: каталоги и
пути}\label{ux433ux434ux435-ux43cux44b-ux43dux430ux445ux43eux434ux438ux43cux441ux44f-ux43aux430ux442ux430ux43bux43eux433ux438-ux438-ux43fux443ux442ux438}}

Введите по очереди:

\begin{verbatim}
pwd
ls
cd ~
pwd
\end{verbatim}

\texttt{pwd} показывает \textbf{текущий каталог} (\emph{working
directory}), а \texttt{ls} --- его содержимое. \texttt{cd} меняет
текущий каталог. Знак \texttt{\textasciitilde{}} означает ваш
\textbf{домашний каталог}; результат \texttt{pwd} после
\texttt{cd\ \textasciitilde{}} будет похож на \texttt{/home/ваше\_имя}.
Это пример, а не одинаковый для всех путь.

В Linux корень дерева каталогов обозначается \texttt{/}. Путь
\texttt{/home/ваше\_имя} начинается с \texttt{/} и называется
\textbf{абсолютным}. Имя \texttt{notes.txt} задаёт путь от текущего
каталога --- это \textbf{относительный} путь. \texttt{.} означает
текущий каталог, \texttt{..} --- родительский. Регистр букв важен:
\texttt{Notes.txt} и \texttt{notes.txt} --- разные имена.

\hypertarget{ux441ux43eux437ux434ux430ux451ux43c-ux443ux447ux435ux431ux43dux44bux439-ux43aux430ux442ux430ux43bux43eux433-ux438-ux444ux430ux439ux43b}{%
\section{Создаём учебный каталог и
файл}\label{ux441ux43eux437ux434ux430ux451ux43c-ux443ux447ux435ux431ux43dux44bux439-ux43aux430ux442ux430ux43bux43eux433-ux438-ux444ux430ux439ux43b}}

Начиная из домашнего каталога, выполните:

\begin{verbatim}
mkdir terminal-practice
cd terminal-practice
pwd
touch notes.txt
ls
\end{verbatim}

\texttt{mkdir} создал каталог; \texttt{cd} перешёл в него.
\texttt{touch} создал пустой файл \texttt{notes.txt} (если файл уже
существует, команда обновит время его изменения). \texttt{ls} теперь
должен показать \texttt{notes.txt}. Если \texttt{mkdir} сообщает
\texttt{File\ exists}, учебная папка уже существует: перейдите в неё
командой \texttt{cd\ terminal-practice} и проверьте содержимое командой
\texttt{ls}, прежде чем продолжать.

Сохраните текст в файле через редактор \textbf{nano}:

\begin{verbatim}
nano notes.txt
\end{verbatim}

Введите в редакторе две строки, например:

\begin{verbatim}
Изучаю терминал Ubuntu.
Мой первый текстовый файл.
\end{verbatim}

\textbf{Сохраните файл в nano}: нажмите \textbf{Ctrl+O} (буква O), затем
\textbf{Enter} для подтверждения имени. Нажмите \textbf{Ctrl+X}, чтобы
выйти. Подсказки внизу nano используют символ \texttt{\^{}} для клавиши
Ctrl: \texttt{\^{}O} означает Ctrl+O. Если при выходе редактор
спрашивает о сохранении несохранённых изменений, подтвердите сохранение
и имя файла.

Убедитесь, что файл содержит ваш текст:

\begin{verbatim}
cat notes.txt
\end{verbatim}

\texttt{cat} выводит содержимое небольшого текстового файла в терминал.
Для длинного файла удобнее \texttt{less\ notes.txt}: прокрутка ---
стрелками, выход --- клавишей \textbf{q}.

\hypertarget{ux43aux43eux43fux438ux440ux443ux435ux43c-ux43fux435ux440ux435ux438ux43cux435ux43dux43eux432ux44bux432ux430ux435ux43c-ux438-ux443ux434ux430ux43bux44fux435ux43c}{%
\section{Копируем, переименовываем и
удаляем}\label{ux43aux43eux43fux438ux440ux443ux435ux43c-ux43fux435ux440ux435ux438ux43cux435ux43dux43eux432ux44bux432ux430ux435ux43c-ux438-ux443ux434ux430ux43bux44fux435ux43c}}

Оставайтесь в \texttt{terminal-practice}. Команды ниже работают только с
файлами учебного примера:

\begin{verbatim}
cp notes.txt notes-copy.txt
ls
mv notes-copy.txt draft.txt
mkdir archive
mv draft.txt archive/
ls archive
rm archive/draft.txt
rmdir archive
ls
\end{verbatim}

\texttt{cp} делает копию, \texttt{mv} переименовывает файл или переносит
его в каталог, \texttt{rm} удаляет файл, а \texttt{rmdir} удаляет
\textbf{пустой} каталог. В конце \texttt{notes.txt} остаётся на месте.
Перед \texttt{rm} проверьте путь через \texttt{ls}: обычно удаление из
терминала не перемещает файл в корзину. Не применяйте эти команды к
незнакомым каталогам.

Из учебного каталога команда \texttt{cd\ ..} вернёт вас на уровень выше
--- домой. Команда \texttt{cd\ \textasciitilde{}} вернёт домой из любого
каталога. Проверяйте результат через \texttt{pwd}.

\hypertarget{ux43aux43eux440ux43eux442ux43aux430ux44f-ux43fux430ux43cux44fux442ux43aux430-ux438-ux43fux43eux43cux43eux449ux44c}{%
\section{Короткая памятка и
помощь}\label{ux43aux43eux440ux43eux442ux43aux430ux44f-ux43fux430ux43cux44fux442ux43aux430-ux438-ux43fux43eux43cux43eux449ux44c}}

\begin{longtable}[]{@{}
  >{\raggedright\arraybackslash}p{(\columnwidth - 2\tabcolsep) * \real{0.3700}}
  >{\raggedright\arraybackslash}p{(\columnwidth - 2\tabcolsep) * \real{0.5900}}@{}}
\toprule
\begin{minipage}[b]{\linewidth}\raggedright
Команда
\end{minipage} & \begin{minipage}[b]{\linewidth}\raggedright
Что делает
\end{minipage} \\
\midrule
\endhead
\texttt{pwd}, \texttt{ls}, \texttt{ls\ -a} & Показать текущий каталог,
его содержимое и скрытые имена. \\
\texttt{cd\ ..}, \texttt{cd\ \textasciitilde{}} & Перейти на уровень
выше или домой. \\
\texttt{mkdir}, \texttt{touch}, \texttt{nano} & Создать каталог, пустой
файл или отредактировать текст. \\
\texttt{cat}, \texttt{less} & Вывести текст или просмотреть его
постранично. \\
\texttt{cp}, \texttt{mv}, \texttt{rm}, \texttt{rmdir} & Копировать,
переместить, удалить файл или пустой каталог. \\
\texttt{grep\ -n\ \textquotesingle{}Ubuntu\textquotesingle{}\ notes.txt}
& Найти строки со словом \texttt{Ubuntu} и показать их номера. \\
\texttt{find\ .\ -name\ \textquotesingle{}*.txt\textquotesingle{}} &
Найти файлы с именем, заканчивающимся на \texttt{.txt}, начиная с
текущего каталога. \\
\texttt{wc\ -l\ notes.txt} & Посчитать строки в файле. \\
\bottomrule
\end{longtable}

В примере с \texttt{find} кавычки вокруг \texttt{*.txt} нужны, чтобы
образец имени обрабатывала команда \texttt{find}, а не оболочка заранее.
Для справки по команде попробуйте \texttt{ls\ -\/-help} или
\texttt{man\ ls}; из справки \texttt{man} выходят клавишей \textbf{q}.
Команды вводите без обратных кавычек --- здесь они лишь выделяют текст
команды.

\hypertarget{ux441ux430ux43cux43eux441ux442ux43eux44fux442ux435ux43bux44cux43dux430ux44f-ux43fux440ux430ux43aux442ux438ux43aux430}{%
\section{Самостоятельная
практика}\label{ux441ux430ux43cux43eux441ux442ux43eux44fux442ux435ux43bux44cux43dux430ux44f-ux43fux440ux430ux43aux442ux438ux43aux430}}

Из домашнего каталога создайте каталог \texttt{my-work}, войдите в него,
создайте файл \texttt{plan.txt} и запишите в нём две собственные строки
через \texttt{nano}. Проверьте содержимое командой \texttt{cat},
создайте копию \texttt{plan-copy.txt}, переименуйте её в
\texttt{backup.txt}, посмотрите список файлов и удалите \textbf{только}
\texttt{backup.txt}. В результате в \texttt{my-work} должен остаться
\texttt{plan.txt} с вашим текстом. Для проверки используйте
\texttt{pwd}, \texttt{ls} и \texttt{cat\ plan.txt}.

Если команда отвечает \texttt{No\ such\ file\ or\ directory}, сначала
проверьте \texttt{pwd} и \texttt{ls}: возможно, вы находитесь не в той
папке либо ошиблись в имени (включая регистр). \texttt{cd} не создаёт
каталог: перед переходом каталог должен существовать. Если команда
отвечает \texttt{command\ not\ found}, проверьте написание команды,
пробелы и раскладку клавиатуры; не повторяйте её с \texttt{sudo} наугад.
Если завис просмотр в \texttt{less} или \texttt{man}, нажмите
\textbf{q}; если нужно прервать выполняющуюся команду, нажмите
\textbf{Ctrl+C}.
